%%%%%%%%%%%%%%%%%%%%%%%%
\chapter{Introdução}\label{intro}
%%%%%%%%%%%%%%%%%%%%%%%%

A crescente adoção de técnicas de aprendizado de máquina em diferentes domínios tem ampliado a necessidade de modelos capazes não apenas de produzir predições acuradas, mas também de oferecer mecanismos compreensíveis para inspeção, validação e interpretação dos resultados. Em áreas nas quais decisões automatizadas podem apoiar diagnósticos, priorização de ações, identificação de padrões ou descoberta de conhecimento, a transparência do modelo passa a ser um critério relevante. Nesse contexto, modelos baseados em regras mantêm importância histórica e metodológica, pois representam o conhecimento aprendido por meio de estruturas simbólicas, geralmente organizadas no formato antecedente-consequente. Essa forma de representação permite que especialistas analisem as condições associadas a determinada decisão e compreendam, ao menos parcialmente, a lógica utilizada pelo modelo.

Essa preocupação é coerente com o debate recente sobre interpretabilidade em aprendizado de máquina. \citeonline{Rudin2019}, por exemplo, é enfática já no título de seu trabalho ao defender: ``Stop explaining black box machine learning models for high stakes decisions and use interpretable models instead''. Embora essa afirmação esteja vinculada a aplicações críticas, ela expressa uma preocupação mais ampla: em muitos cenários, explicar uma decisão pode ser tão importante quanto acertá-la. De forma complementar, \citeonline{molnar2022} discute que a interpretabilidade contribui para confiança, auditabilidade e compreensão do comportamento dos modelos. No caso específico de modelos baseados em regras, \citeonline{margot2021interpretability} destacam que a interpretabilidade pode ser analisada a partir de propriedades estruturais do modelo, como quantidade de regras, tamanho das regras e complexidade da representação.

Entre os modelos baseados em regras, a classificação associativa (CA) destaca-se por combinar mineração de regras de associação e classificação supervisionada. A proposta seminal de \citeonline{liu1998}, cujo título explicita a ideia de ``Integrating Classification and Association Rule Mining'', estabeleceu uma linha de investigação na qual regras de associação são adaptadas para fins preditivos. Nesse contexto, utilizam-se regras de associação classificativas (RACs), geralmente expressas na forma $A \Rightarrow c$, em que $A$ representa um conjunto de itens no antecedente e $c$ corresponde à classe do problema. Assim, diferentemente da mineração de regras de associação tradicional, cujo objetivo é descobrir relações frequentes entre itens, a CA direciona a extração de regras para a construção de modelos capazes de classificar novas instâncias.

A CA tornou-se uma abordagem relevante porque concilia capacidade preditiva e representação interpretável do conhecimento. Algoritmos como CBA, CMAR, MCAR e CPAR consolidaram diferentes estratégias para geração, ordenação, poda e aplicação de regras classificativas \cite{liu1998,Li2001,1387030,Dehkordi2006}. Além disso, revisões e bibliotecas recentes demonstram que a área continua ativa e aplicável a diferentes domínios \cite{THABTAH:2007:RAC,Hahsler2019254,Padillo2020}. Segundo \citeonline{Padillo2020}, a classificação associativa apresenta vantagens relevantes em relação a outras abordagens interpretáveis, especialmente por permitir que o modelo seja representado por regras explícitas, passíveis de análise por especialistas.

Apesar dessas vantagens, a interpretabilidade dos classificadores associativos não é uma consequência automática do uso de regras. Um modelo baseado em regras pode ser formalmente interpretável, mas praticamente difícil de compreender quando contém muitas regras, regras extensas, antecedentes redundantes ou padrões pouco relevantes. Assim, a qualidade de um classificador associativo depende não apenas de sua acurácia, mas também da compacidade e da utilidade do conjunto de regras induzido. Em outras palavras, uma regra isolada pode ser simples de ler; milhares delas, não. A tradição dos modelos simbólicos ajuda, mas não faz milagre estatístico.

Esse problema torna-se mais evidente quando classificadores associativos são aplicados a conjuntos de dados de alta dimensionalidade. A mineração de regras de associação depende da geração de itemsets frequentes, etapa reconhecidamente custosa do ponto de vista computacional. Conforme proposto nos trabalhos clássicos de \citeonline{agrawal1993mining} e \citeonline{apriori}, a extração de regras envolve identificar combinações de itens que satisfazem limiares mínimos de suporte e, posteriormente, gerar regras que atendam a critérios como confiança. Entretanto, à medida que o número de itens aumenta, o espaço de busca cresce de maneira combinatória, tornando a geração de itemsets frequentes a etapa mais dispendiosa do processo.

Esse crescimento é particularmente crítico em CA porque os atributos originais do conjunto de dados são convertidos em itens no formato \texttt{<atributo-valor>}. Desse modo, um único atributo categórico pode originar múltiplos itens, a depender da cardinalidade de seus valores. Além disso, atributos contínuos frequentemente precisam ser discretizados antes da mineração de regras, transformando intervalos numéricos em itens categóricos. A discretização, embora necessária para muitos algoritmos associativos, pode ampliar ainda mais o conjunto de itens quando vários atributos contínuos são divididos em múltiplos intervalos \cite{conf/ijcai/FayyadI93}. Consequentemente, o número de combinações possíveis aumenta, elevando o custo computacional e favorecendo a geração de um conjunto excessivo de regras candidatas.

A alta dimensionalidade afeta a CA em diferentes níveis. No primeiro nível, amplia o espaço de itens disponíveis para a mineração. No segundo, aumenta o número de itemsets candidatos que precisam ser avaliados. No terceiro, potencializa a quantidade de regras geradas a partir desses itemsets. No quarto, dificulta a etapa de ordenação e poda, pois há mais regras a serem comparadas, avaliadas e descartadas. Por fim, no quinto nível, compromete a interpretabilidade do modelo final, já que um número elevado de regras reduz a capacidade de análise humana. Portanto, a dimensionalidade não afeta apenas a entrada do classificador; ela se propaga por todo o pipeline associativo.

Além do impacto computacional, a alta dimensionalidade também pode prejudicar a qualidade das regras. A presença de atributos irrelevantes, redundantes ou ruidosos aumenta a probabilidade de geração de padrões espúrios, isto é, regras que apresentam bom comportamento no conjunto de treinamento, mas baixa capacidade de generalização. Em bases com muitos atributos, algumas combinações podem parecer relevantes apenas por acaso, especialmente quando os dados são esparsos ou quando há grande número de valores possíveis. Esse fenômeno é compatível com a discussão mais ampla sobre a maldição da dimensionalidade, na qual o aumento do número de variáveis pode tornar o espaço de dados mais esparso, dificultar a identificação de padrões robustos e elevar o risco de sobreajuste \cite{Dhal2022,Theng2024,BARBIERI2024123667}.

No caso dos classificadores associativos, esse problema assume contornos específicos. Enquanto em muitos classificadores a alta dimensionalidade afeta principalmente a estimação dos parâmetros ou a separabilidade das classes, em CA ela interfere diretamente na etapa de descoberta de padrões. Como as regras são formadas por combinações de itens, atributos irrelevantes não são apenas variáveis sem utilidade; eles passam a compor combinações candidatas, aumentando o número de itemsets e regras avaliadas. Assim, um atributo pouco informativo pode contribuir para o aumento do custo de mineração e para a geração de regras redundantes ou pouco interpretáveis. A consequência prática é dupla: o modelo pode se tornar mais caro de induzir e mais difícil de compreender.

Nesse cenário, a seleção de atributos, ou \textit{feature selection} (FS), surge como uma estratégia natural para mitigar os efeitos da alta dimensionalidade. De acordo com \citeonline{Dhal2022}, a FS é uma etapa relevante em aprendizado de máquina por permitir a remoção de atributos irrelevantes, redundantes ou ruidosos. \citeonline{Theng2024}, em uma revisão abrangente sobre mais de duas décadas de pesquisa em seleção de atributos, destacam que a FS busca reduzir a dimensionalidade preservando, ou eventualmente melhorando, o desempenho dos modelos. Em termos gerais, a FS procura identificar um subconjunto de atributos suficientemente informativo para representar o problema, evitando que todos os atributos originais sejam utilizados de maneira indiscriminada.

A aplicação de FS em CA possui uma motivação ainda mais forte do que em classificadores tradicionais. Ao reduzir o número de atributos antes da mineração, reduz-se também o número de itens derivados desses atributos. Como consequência, diminui-se o espaço de busca para geração de itemsets e regras. Essa redução pode impactar positivamente o tempo de execução, o consumo de memória, a quantidade de regras geradas e a interpretabilidade do classificador final. Portanto, em CA, a FS não atua apenas como uma técnica de pré-processamento voltada à melhoria de desempenho preditivo; ela funciona como um mecanismo de controle estrutural do processo de indução.

Os métodos de FS são geralmente organizados em três grandes abordagens: filtros, wrappers e embedded \cite{Dhal2022,Theng2024}. Métodos filtro avaliam os atributos por meio de critérios estatísticos, informacionais, de similaridade ou de dependência, independentemente do classificador utilizado. Métodos wrapper avaliam subconjuntos de atributos com base no desempenho de um classificador, enquanto métodos embedded incorporam a seleção ao próprio processo de treinamento. No contexto deste trabalho, os métodos filtro são particularmente relevantes por apresentarem menor custo computacional, independência em relação ao classificador e maior adequação a cenários de alta dimensionalidade \cite{SolorioFernandez2020,Bouchlaghem2022,Theng2024}.

Apesar de suas vantagens, a FS isolada não resolve integralmente os desafios da CA. O primeiro motivo é que diferentes métodos de FS incorporam diferentes noções de relevância. Um método baseado em ganho de informação pode priorizar atributos com maior redução de entropia em relação à classe; um método baseado em correlação pode favorecer relações lineares ou monotônicas; um método baseado em vizinhança, como ReliefF, pode capturar diferenças locais entre instâncias; e métodos multivariados podem considerar redundância ou dependência entre atributos. Assim, a escolha de um único seletor introduz um viés metodológico, pois o subconjunto final passa a refletir a preferência estatística daquele critério específico \cite{KHAIRE20221060,BARBIERI2024123667}.

O segundo motivo é que muitos métodos filtro, sobretudo os univariados, avaliam atributos de maneira isolada. Essa característica é eficiente, mas pode ser limitada em problemas nos quais a capacidade discriminativa emerge da combinação entre atributos. Em CA, essa limitação é especialmente importante, pois o modelo é justamente construído a partir de combinações de itens no antecedente das regras. Um atributo pode apresentar baixa relevância individual, mas contribuir para uma regra útil quando combinado com outros atributos. Assim, a remoção prematura de atributos com base apenas em relevância individual pode eliminar informações importantes para a formação de RACs.

O terceiro motivo está relacionado à redundância. Métodos de FS podem reduzir dimensionalidade, mas nem sempre eliminam adequadamente atributos redundantes. Dois atributos podem apresentar alta relevância individual em relação à classe e, ainda assim, carregar informação semelhante. Em classificadores associativos, atributos redundantes podem originar regras redundantes, com antecedentes semelhantes e consequentes idênticos. Isso aumenta o tamanho do modelo sem necessariamente agregar novo conhecimento. Portanto, selecionar atributos relevantes não é suficiente; é necessário considerar também a redundância e a complementaridade entre eles \cite{KHAIRE20221060,Theng2024}.

O quarto motivo refere-se à definição do ponto de corte. Muitos métodos de FS produzem rankings de atributos, mas não determinam automaticamente quantos atributos devem ser mantidos. A seleção dos $k$ melhores atributos, de um percentual fixo ou de todos os atributos acima de um determinado limiar é uma decisão metodológica sensível. Um corte permissivo pode preservar atributos irrelevantes e manter o espaço de busca excessivamente grande. Um corte restritivo pode eliminar atributos úteis e prejudicar a qualidade das regras. Em CA, essa decisão é ainda mais crítica, pois o número de atributos selecionados afeta diretamente a quantidade de itens, itemsets e regras. Como discutem \citeonline{seijopardo2019automatic}, já no título do trabalho ``On developing an automatic threshold applied to feature selection ensembles'', a definição automática de limiares em seleção de atributos é um problema próprio, não apenas um detalhe operacional.

O quinto motivo é que a FS externa pode ficar dissociada do comportamento do classificador associativo. Quando a seleção é aplicada antes da indução, o seletor avalia atributos sem considerar diretamente como eles participarão da geração, ordenação e poda de regras. Um atributo selecionado por uma métrica estatística pode não contribuir para regras úteis; inversamente, um atributo descartado pode ser importante em combinações específicas. Essa dissociação revela uma tensão metodológica: reduzir atributos antes da mineração simplifica o processo, mas pode ignorar a lógica combinatória que caracteriza os modelos associativos.

Essas limitações indicam que a FS deve ser analisada com cautela no contexto da CA. Aplicá-la mecanicamente como etapa de pré-processamento pode reduzir dimensionalidade, mas não garante melhoria de desempenho, redução efetiva de regras ou aumento de interpretabilidade. O efeito real da FS depende do método utilizado, do tipo dos atributos, da discretização, dos limiares de suporte e confiança, do algoritmo associativo, do critério de ordenação das regras e da estratégia de poda. Portanto, avaliar FS em CA exige observar não apenas métricas preditivas, mas também medidas estruturais do modelo, como número de regras, tamanho das regras, tempo de execução, cobertura e estabilidade dos atributos selecionados.

Diante dessas limitações, o \textit{ensemble de feature selection} (EFS) apresenta-se como uma alternativa promissora. Em vez de depender de um único seletor, o EFS combina múltiplas evidências de relevância para produzir uma decisão agregada sobre os atributos. \citeonline{bolon2018scalability} discutem que ensembles de seleção de atributos podem contribuir para melhorar a robustez e a estabilidade dos subconjuntos selecionados. De modo semelhante, \citeonline{mera2021framework} tratam a combinação de métodos de FS como um problema de consenso, no qual diferentes seletores funcionam como fontes complementares de informação. Essa lógica desloca a decisão do ``melhor seletor isolado'' para a construção de uma decisão coletiva sobre a importância dos atributos.

A principal justificativa para o uso de EFS está na diversidade dos seletores. Como diferentes métodos capturam diferentes aspectos dos dados, a combinação de seus resultados pode reduzir a dependência de um critério específico. Em um conjunto de dados, determinado atributo pode ser altamente relevante segundo uma medida informacional, moderadamente relevante segundo uma medida estatística e pouco relevante segundo uma medida de distância. Um ensemble pode agregar essas perspectivas, produzindo um ranking ou subconjunto mais equilibrado. Em CA, essa robustez é desejável porque pequenas alterações no conjunto de atributos podem alterar significativamente o conjunto de regras geradas.

Além da robustez, o EFS pode contribuir para a estabilidade. A estabilidade de FS refere-se à consistência dos atributos selecionados quando há variações nos dados, nas amostras ou nos particionamentos experimentais \cite{KHAIRE20221060,BARBIERI2024123667}. Em problemas interpretáveis, essa propriedade é importante porque a seleção de atributos não deve apenas produzir bom desempenho; ela deve também indicar, de forma consistente, quais variáveis são relevantes. Em classificadores associativos, a instabilidade dos atributos selecionados pode gerar instabilidade nas regras, dificultando a interpretação do conhecimento extraído. Assim, um EFS mais estável pode contribuir para modelos associativos mais confiáveis.

O EFS também pode favorecer a complementaridade entre métodos univariados e multivariados. Métodos univariados são eficientes e escaláveis, mas podem ignorar interações. Métodos multivariados tendem a capturar redundância e dependência, mas podem ser mais custosos. Ao combinar seletores distintos, torna-se possível explorar diferentes critérios sem depender integralmente de um único tipo de avaliação. Estratégias baseadas em agregação de rankings, votação, métodos eleitorais, tomada de decisão multicritério e otimização multiobjetivo têm sido investigadas com esse objetivo \cite{drotar2019ensemble,hashemi2021pareto,abellana2022ensemble,kaur2023mcdm,abasabadi2021automatic}.

Entretanto, o EFS também não deve ser visto como solução automática. Ele introduz novas decisões metodológicas: quais seletores combinar, como normalizar rankings ou scores, como tratar empates, como ponderar seletores, como agregar evidências e como definir o subconjunto final de atributos. Além disso, a etapa de \textit{threshold} permanece central. Após a agregação, ainda é necessário determinar quais atributos serão efetivamente preservados. Assim, o EFS pode reduzir a dependência de um único seletor, mas não elimina a necessidade de escolhas experimentais rigorosas. Em termos simples: trocar um seletor por um comitê melhora a conversa, mas ainda é preciso decidir quem entra na sala.

A literatura analisada neste trabalho indica que há esforços relevantes para integrar FS e CA, mas também evidencia lacunas importantes. Estudos como os de \citeonline{shaharanee2015irrelevant}, \citeonline{sowan2023hybrid}, \citeonline{wu2024novel}, \citeonline{nandhini2020performance}, \citeonline{gayathri2015efficient}, \citeonline{cicek2021assessment}, \citeonline{alhawari2021classification}, \citeonline{ashi2021enhancing}, \citeonline{alazaidah2021feature} e \citeonline{piri2019exploring} demonstram diferentes formas de empregar FS em classificadores associativos. Essas abordagens incluem seleção externa de atributos, seleção dinâmica durante a indução de regras, ponderação de itens, uso de métodos embedded, wrappers, filtros e estratégias híbridas.

Apesar dessa diversidade, observa-se que a maior parte dos estudos utiliza a FS como etapa externa ao processo de indução, isto é, como pré-processamento. Essa estratégia é simples e modular, mas nem sempre explora a relação entre seleção de atributos e geração de regras. Além disso, muitos trabalhos avaliam predominantemente o desempenho preditivo, enquanto aspectos diretamente associados à natureza dos classificadores associativos, como número de regras, redução de itemsets, tamanho do modelo, estabilidade e interpretabilidade, recebem menor atenção. Essa limitação é relevante porque a contribuição esperada da FS em CA não deve ser medida apenas pela acurácia, mas também pela capacidade de simplificar o modelo e tornar o conhecimento extraído mais compreensível.

Outra lacuna refere-se ao uso ainda limitado de estratégias de EFS especificamente orientadas à classificação associativa. Embora o EFS seja amplamente discutido na literatura geral de seleção de atributos, sua aplicação em CA ainda carece de investigação sistemática. Em particular, permanece em aberto como combinar diferentes métodos filtro para reduzir o espaço de busca associativo, como definir o \textit{threshold} após a agregação dos rankings e como avaliar o impacto dessa seleção sobre o conjunto final de RACs. Essa lacuna é relevante porque a seleção de atributos em CA não afeta apenas o desempenho do classificador; ela modifica o próprio processo de descoberta de regras.

Assim, a lacuna central investigada neste trabalho consiste na ausência de uma análise sistemática sobre estratégias de seleção de atributos para classificadores associativos que considere simultaneamente redução dimensional, desempenho preditivo, complexidade computacional, estabilidade e interpretabilidade. A literatura mostra que FS pode ser útil para CAs, mas ainda há espaço para investigar como diferentes seletores e estratégias de ensemble afetam a geração de itemsets, o número de regras e a qualidade estrutural do modelo. Em especial, falta compreender se a combinação de múltiplos critérios de seleção pode produzir subconjuntos de atributos mais robustos do que aqueles obtidos por métodos isolados.

Diante desse contexto, este trabalho investiga estratégias de seleção de atributos aplicadas à classificação associativa, com foco em métodos filtro e ensembles de seleção de atributos. A proposta parte da premissa de que a redução adequada do conjunto de atributos pode diminuir o espaço de busca, reduzir a quantidade de regras geradas e melhorar a interpretabilidade dos modelos, desde que o processo de seleção preserve informações relevantes para a tarefa de classificação. Além disso, considera-se que a combinação de diferentes métodos de FS pode mitigar limitações individuais dos seletores, produzindo subconjuntos mais estáveis e menos dependentes de um único critério de relevância.

A contribuição esperada deste estudo não se limita a comparar acurácias entre modelos com e sem FS. O foco está em compreender como a seleção de atributos altera o comportamento do classificador associativo ao longo de seu pipeline. Isso inclui observar a redução de atributos, o impacto sobre a mineração de itemsets, a quantidade de regras extraídas, a estrutura do classificador final, o tempo de execução e a estabilidade dos subconjuntos selecionados. Essa perspectiva é necessária porque, em classificadores associativos, desempenho e interpretabilidade são dimensões complementares. Um modelo ligeiramente mais acurado, mas composto por um número excessivo de regras, pode ser menos útil do que um modelo compacto, estável e suficientemente preciso.

Portanto, este trabalho posiciona a seleção de atributos como um componente metodológico central para o avanço da classificação associativa em cenários de alta dimensionalidade. A investigação proposta busca contribuir para a compreensão de quando e como a FS pode beneficiar CAs, quais limitações permanecem quando métodos isolados são utilizados e em que medida estratégias de ensemble podem produzir seleções mais robustas. Ao articular redução dimensional, geração de regras e interpretabilidade, a pesquisa procura avançar em uma direção ainda pouco consolidada na literatura: tratar a FS não apenas como etapa anterior ao classificador, mas como fator determinante para a qualidade do conhecimento associativo produzido.

Os capítulos seguintes apresentam os fundamentos conceituais e a revisão sistemática que sustentam essa investigação. Inicialmente, são discutidos os conceitos de regras de associação, medidas objetivas, classificação associativa, seleção de atributos, métodos filtro, ensemble de feature selection, estratégias de \textit{threshold} e critérios de avaliação. Em seguida, é apresentada a revisão sistemática da literatura, conduzida para identificar como a FS tem sido incorporada aos classificadores associativos, quais estratégias têm sido utilizadas e quais lacunas permanecem em aberto. Essa organização permite estabelecer a base teórica e empírica necessária para a continuidade da pesquisa.
